% Narrativa generada automaticamente desde clustering_participantes (M2): NO editar a mano.
% Cada cifra esta interpolada desde compute_clustering; regenerar actualiza el texto.
Sobre 5.815 participantes de pregrado con datos completos (de 5.821 participantes de pregrado y 8.036 participantes en todos los niveles), el agrupamiento K-means identifica 2 perfiles internos a partir de cinco rasgos académicos y demográficos estandarizados. En concreto, el grupo «Grupo en percentiles altos» reúne 3.761 estudiantes (64,7\%), caracterizado sobre todo por percentil de pga (programa); el grupo «Grupo en percentiles bajos» reúne 2.054 estudiantes (35,3\%), caracterizado sobre todo por percentil de pga (programa). Entre las variables de perfilado externas, «ICFES puntaje total» es la que más diferencia a los grupos; sus medias se calculan solo sobre quienes tienen dato (cobertura parcial en los puntajes ICFES y en el estrato). Estos agrupamientos son descriptivos, no categorías administrativas ni relaciones causales; la media de estrato es un indicador ordinal orientativo y todas las medias de cobertura parcial deben leerse junto a su porcentaje de dato disponible.
